\documentclass[12pt]{article}
\usepackage[utf8]{inputenc}
\usepackage[T1]{fontenc}
\usepackage[margin=1in]{geometry}
\usepackage{setspace}
\usepackage{xcolor}
\usepackage[colorlinks=true, urlcolor=blue!60!black, linkcolor=blue!60!black]{hyperref}
\setlength{\parindent}{2em}
\setlength{\parskip}{1.25em}
\renewcommand{\baselinestretch}{1.0}
\begin{document}
\begin{center}
{\LARGE \textit{Roe v. Wade}, 410 U.S. 113 (1973)\par}
\vspace{0.8em}
{\large Case Brief\par}
\end{center}
\vspace{1em}

\section{Memory Jogger}

In \href{https://www.courtlistener.com/opinion/108713/roe-v-wade/}{\textit{Roe v. Wade}}, the Supreme Court held that the Fourteenth Amendment's Due Process Clause protects a fundamental right of personal privacy broad enough to encompass a woman's decision to terminate a pregnancy, governed by a trimester framework balancing maternal health and fetal viability (later modified in \textit{Planned Parenthood of Southeastern Pennsylvania v. Casey} (1992) and overruled in \href{https://www.supremecourt.gov/opinions/21pdf/19-1392_6j37.pdf}{\textit{Dobbs v. Jackson Women's Health Organization}} (2022)).

\section{Detailed Case Facts}

In March 1970, Norma McCorvey---suing under the pseudonym ``Jane Roe''---was an unmarried, pregnant resident of Dallas County, Texas, who wanted to terminate her pregnancy through a safe, clinical procedure performed by a licensed physician (\href{https://supreme.justia.com/cases/federal/us/410/113/}{410 U.S. at 120}; \href{https://www.oyez.org/cases/1971/70-18}{Oyez}). Under Articles 1191--1194 and 1196 of the Texas Penal Code, first enacted in 1854 and substantially unchanged since 1857, procuring or attempting an abortion was a felony carrying two to five years' imprisonment unless ``procured or attempted by medical advice for the purpose of saving the life of the mother'' (\href{https://supreme.justia.com/cases/federal/us/410/113/}{410 U.S. at 117--19 \& n.1}). Because continuing her pregnancy did not threaten her life, Roe could not obtain a legal abortion in Texas and could not afford to travel to a state where abortion was lawful (\href{https://supreme.justia.com/cases/federal/us/410/113/}{410 U.S. at 120}).

Represented by attorneys Linda Coffee and Sarah Weddington, Roe filed a federal class action against Henry Wade, the District Attorney of Dallas County, seeking declaratory and injunctive relief on behalf of herself and all similarly situated women (\href{https://supreme.justia.com/cases/federal/us/410/113/}{410 U.S. at 120}). Two additional sets of plaintiffs joined the consolidated litigation:
\begin{itemize}
\item \textbf{Dr. James Hubert Hallford}, a Texas-licensed physician who intervened as a plaintiff-intervenor, alleging that he had been previously arrested under the Texas abortion statutes and currently faced two pending criminal prosecutions in Dallas County, and contending that Article 1196's exception was unconstitutionally vague (\href{https://supreme.justia.com/cases/federal/us/410/113/}{410 U.S. at 120--21, 125--26}).
\item \textbf{John and Mary Doe}, a childless, non-pregnant married couple who filed a companion complaint alleging that Mrs. Doe suffered from a ``neural-chemical'' disorder, had been medically advised to avoid pregnancy, had stopped taking birth control pills on medical advice, and feared future contraceptive failure and an inability to obtain a legal abortion in Texas (\href{https://supreme.justia.com/cases/federal/us/410/113/}{410 U.S. at 121}).
\end{itemize}

Roe gave birth before the litigation reached final resolution, prompting Texas to argue that her claim had become moot (\href{https://supreme.justia.com/cases/federal/us/410/113/}{410 U.S. at 124--25}). The Supreme Court heard \textit{Roe} alongside its companion case, \href{https://supreme.justia.com/cases/federal/us/410/179/}{\textit{Doe v. Bolton}}, 410 U.S. 179 (1973), which challenged Georgia's more modern, American Law Institute-modeled abortion statute.

\section{Procedural History}

\begin{enumerate}
\item \textbf{U.S. District Court for the Northern District of Texas (\href{https://law.justia.com/cases/federal/district-courts/FSupp/314/1217/1472349/}{314 F. Supp. 1217 (N.D. Tex. 1970)}):} A three-judge federal district court consolidated the Roe and Doe actions and heard motions to dismiss and for summary judgment on May 22, 1970. On June 17, 1970, the court held that Roe and Dr. Hallford had standing and presented justiciable controversies, while dismissing the Does' complaint as too speculative. On the merits, the district court declared Articles 1191--1194 and 1196 facially unconstitutional as vague and as an overbroad infringement of the ``fundamental right of single women and married persons to choose whether to have children'' protected by the Ninth Amendment through the Fourteenth Amendment. Citing federal abstention principles, however, the court denied injunctive relief (\href{https://supreme.justia.com/cases/federal/us/410/113/}{410 U.S. at 122}).
\item \textbf{Direct Appeal to the U.S. Supreme Court:} Pursuant to 28 U.S.C. \S~1253, Roe, the Does, and Dr. Hallford appealed the denial of injunctive relief directly to the Supreme Court, and Wade cross-appealed the grant of declaratory relief (both sides also filed protective appeals in the U.S. Court of Appeals for the Fifth Circuit, which held them in abeyance) (\href{https://supreme.justia.com/cases/federal/us/410/113/}{410 U.S. at 122}).
\item \textbf{Argument, Reargument, and Decision:} The Supreme Court postponed jurisdiction to the merits (402 U.S. 941 (1971)), heard initial oral argument on December 13, 1971 (before a seven-Justice bench following the retirements of Justices Black and Harlan), ordered reargument on October 11, 1972 (after Justices Powell and Rehnquist joined the Court), and issued its decision on January 22, 1973, affirming in part and reversing in part (\href{https://supreme.justia.com/cases/federal/us/410/113/}{410 U.S. at 113, 115, 166--67}; \href{https://www.oyez.org/cases/1971/70-18}{Oyez}).
\end{enumerate}

\section{Judicial Votes}

The decision was \textbf{7--2} to invalidate the Texas criminal abortion statutes (\href{https://supreme.justia.com/cases/federal/us/410/113/}{410 U.S. at 115}):
\begin{itemize}
\item \textbf{Majority (7):} Justice Harry A. Blackmun (author of the opinion of the Court), joined by Chief Justice Warren E. Burger and Justices William O. Douglas, William J. Brennan Jr., Potter Stewart, Thurgood Marshall, and Lewis F. Powell Jr.
\item \textbf{Concurring Opinions (3):} Chief Justice Burger (published at \href{https://supreme.justia.com/cases/federal/us/410/179/}{410 U.S. at 207}, applying to both \textit{Roe} and \textit{Doe v. Bolton}), Justice Douglas (published at \href{https://supreme.justia.com/cases/federal/us/410/179/}{410 U.S. at 209}, applying to both cases, though disagreeing in a footnote with dismissing Dr. Hallford's complaint under \textit{Younger v. Harris}), and Justice Stewart (\href{https://supreme.justia.com/cases/federal/us/410/113/}{410 U.S. at 167}).
\item \textbf{Dissenting Opinions (2):} Justice Byron R. White, joined by Justice William H. Rehnquist (published at \href{https://supreme.justia.com/cases/federal/us/410/179/}{410 U.S. at 221}, applying to both \textit{Roe} and \textit{Doe}); and Justice Rehnquist separately (\href{https://supreme.justia.com/cases/federal/us/410/113/}{410 U.S. at 171}).
\end{itemize}

\section{Holding}

\textbf{Judgment:} Affirmed in part and reversed in part. The Court affirmed the district court's declaratory judgment striking down the Texas criminal abortion statutes (Articles 1191--1194 and 1196) as a unit, affirmed the dismissal of the Does' complaint for lack of standing, reversed the grant of declaratory relief to Dr. Hallford and ordered his complaint in intervention dismissed under \textit{Younger v. Harris} and \textit{Samuels v. Mackell}, and found it unnecessary to decide whether injunctive relief should have issued (\href{https://supreme.justia.com/cases/federal/us/410/113/}{410 U.S. at 115, 127, 129, 166--67}).

\textbf{Rule Announced:}
\begin{enumerate}
\item \textbf{Justiciability:} Pregnancy is ``capable of repetition, yet evading review,'' so the termination of a plaintiff's pregnancy before appellate review does not moot her challenge (\href{https://supreme.justia.com/cases/federal/us/410/113/}{410 U.S. at 125}).
\item \textbf{Substantive Due Process Right:} The Due Process Clause of the Fourteenth Amendment protects a fundamental right of personal privacy that encompasses a woman's qualified decision to terminate her pregnancy, subject to strict scrutiny (requiring a compelling state interest and narrowly drawn regulation). A fetus is not a ``person'' within the meaning of the Fourteenth Amendment (\href{https://supreme.justia.com/cases/federal/us/410/113/}{410 U.S. at 153--56, 158}).
\item \textbf{Trimester Framework:} Because the State's legitimate interests in maternal health and the ``potentiality of human life'' grow as pregnancy progresses, regulation is governed by three stages (\href{https://supreme.justia.com/cases/federal/us/410/113/}{410 U.S. at 162--65}):
  \begin{itemize}
  \item \textit{First Trimester:} Prior to approximately the end of the first trimester, the abortion decision and its effectuation must be left to the medical judgment of the pregnant woman's attending physician, free of state regulation.
  \item \textit{Second Trimester:} Subsequent to approximately the end of the first trimester and prior to viability, the State may regulate the abortion procedure in ways reasonably related to maternal health (such as provider qualifications and facility licensing).
  \item \textit{Post-Viability (Third Trimester):} Subsequent to fetal viability (typically 24--28 weeks), the State's interest in potential life becomes compelling, and it may regulate or even proscribe abortion except where necessary, in appropriate medical judgment, to preserve the life or health of the mother.
  \end{itemize}
\item \textbf{Physician Licensing:} A State may define ``physician'' as a state-licensed physician and proscribe abortions performed by non-physicians (\href{https://supreme.justia.com/cases/federal/us/410/113/}{410 U.S. at 165}).
\end{enumerate}

\textbf{Holding vs. Dicta:} The core holding invalidated Texas's sweeping ban that allowed abortion only to save the mother's life regardless of pregnancy stage, established Roe's non-mootness and the dismissal of Hallford and the Does, and held that a fetus is not a Fourteenth Amendment ``person.'' By contrast, extensive historical surveys (ancient Greek and Roman attitudes, the Pythagorean origins of the Hippocratic Oath, and medieval ``mediate animation''), the Court's illustrations of permissible second-trimester health regulations (facility licensing, after-care), and footnote 67's reservation of paternal and parental consent issues operated as contextual dicta or explicit reservations rather than necessary adjudications of Texas's single-exception statute (\href{https://supreme.justia.com/cases/federal/us/410/113/}{410 U.S. at 130--47, 163, 165 n.67}). (Note: \href{https://www.law.cornell.edu/gender-justice/resource/roe_v._wade}{\textit{Planned Parenthood v. Casey}}, 505 U.S. 833 (1992), later discarded the trimester framework in favor of the ``undue burden'' test, and \href{https://www.supremecourt.gov/opinions/21pdf/19-1392_6j37.pdf}{\textit{Dobbs v. Jackson Women's Health Organization}}, 597 U.S. 215 (2022), overruled both \textit{Roe} and \textit{Casey}.)

\section{Analysis of Opinions}

\subsection{Majority Opinion (Justice Blackmun)}

Justice Blackmun first resolved jurisdiction under 28 U.S.C. \S~1253 and threshold justiciability questions. Because normal human gestation (266 days) is shorter than the appellate process, pregnancy presents a classic ``capable of repetition, yet evading review'' exception to mootness (\href{https://supreme.justia.com/cases/federal/us/410/113/}{410 U.S. at 125}). Dr. Hallford's intervention, however, was barred by \textit{Samuels v. Mackell}, 401 U.S. 66 (1971), and \textit{Younger v. Harris}, 401 U.S. 37 (1971), because he could raise his constitutional defenses in his pending state criminal prosecutions (\href{https://supreme.justia.com/cases/federal/us/410/113/}{410 U.S. at 126--27}). The Does' claim rested on a chain of speculative contingencies and failed Article III's case-or-controversy requirement (\href{https://supreme.justia.com/cases/federal/us/410/113/}{410 U.S. at 128--29}).

On the merits, Blackmun conducted a lengthy historical survey---spanning ancient practice, the Hippocratic Oath, English common law's ``quickening'' distinction, nineteenth-century American statutes, and positions of the AMA, APHA, and ABA---to conclude that restrictive criminal abortion bans were of relatively recent nineteenth-century vintage and that women had enjoyed broader freedom to terminate early pregnancies at common law and when the Constitution was adopted (\href{https://supreme.justia.com/cases/federal/us/410/113/}{410 U.S. at 129--47}). Evaluating three historical justifications for abortion bans, he dismissed Victorian sexual morality as unasserted by Texas, found that modern antisepsis and medical techniques made first-trimester abortion at least as safe as childbirth (weakening early-stage maternal-health bans), and recognized the State's legitimate interest in prenatal life (\href{https://supreme.justia.com/cases/federal/us/410/113/}{410 U.S. at 148--52}).

Turning to constitutional doctrine, Blackmun acknowledged that the Constitution does not explicitly mention a right of privacy, but synthesized prior cases---including \textit{Meyer v. Nebraska}, \textit{Pierce v. Society of Sisters}, \textit{Skinner v. Oklahoma}, \textit{Griswold v. Connecticut}, \textit{Loving v. Virginia}, and \textit{Eisenstadt v. Baird}---to locate a fundamental right of personal privacy in the Fourteenth Amendment's concept of personal liberty (\href{https://supreme.justia.com/cases/federal/us/410/113/}{410 U.S. at 152--53}). Because denying the abortion choice inflicts medical, psychological, financial, and social detriments, the right encompasses the abortion decision, though it is not absolute (\href{https://supreme.justia.com/cases/federal/us/410/113/}{410 U.S. at 153--54}).

Blackmun rejected Texas's argument that a fetus is a ``person'' under the Fourteenth Amendment, demonstrating that nearly every constitutional use of ``person'' (such as citizenship, apportionment, and officeholding qualifications) applies only postnatally, and noting that Texas's own statute did not treat abortion as murder or punish the woman as an accomplice (\href{https://supreme.justia.com/cases/federal/us/410/113/}{410 U.S. at 156--58 \& n.54}). Declining to resolve when life begins amid philosophical and theological disagreement, the Court held that Texas could not override a woman's fundamental right by adopting one theory of life, and instead balanced the competing interests through the three-stage trimester framework (\href{https://supreme.justia.com/cases/federal/us/410/113/}{410 U.S. at 159--65}). Because Texas's Article 1196 made no distinction by pregnancy stage and allowed only life-saving abortions, the Texas statutes fell as a unit (\href{https://supreme.justia.com/cases/federal/us/410/113/}{410 U.S. at 164, 166}).

\subsection{Concurring Opinion (Chief Justice Burger)}

Writing for both \textit{Roe} and \textit{Doe}, Chief Justice Burger agreed that the Texas and Georgia statutes impermissibly restricted abortions needed to protect maternal health broadly defined, and noted that women pregnant from rape or incest should not have to depend on uncodified prosecutorial discretion (\href{https://supreme.justia.com/cases/federal/us/410/179/}{410 U.S. at 207--08}). While expressing slight unease at the Court's reliance on scientific and medical data and stating he would have permitted a two-physician certification requirement, he emphasized that the Court's rulings did not establish ``abortions on demand'' and trusted physicians' professional standards (\href{https://supreme.justia.com/cases/federal/us/410/179/}{410 U.S. at 208}).

\subsection{Concurring Opinion (Justice Douglas)}

Also writing for both cases, Justice Douglas clarified the doctrinal foundation of privacy, explaining that while the Ninth Amendment does not itself create federally enforceable rights, the ``Blessings of Liberty'' retained by the people inform the Fourteenth Amendment's ``liberty'' (\href{https://supreme.justia.com/cases/federal/us/410/179/}{410 U.S. at 210--11}). He organized privacy into three tiers: (1) absolute First Amendment autonomy over intellect and personality; (2) fundamental freedom of choice in marriage, divorce, procreation, contraception, and childrearing (subject to strict scrutiny); and (3) fundamental freedom to care for one's health and person (\href{https://supreme.justia.com/cases/federal/us/410/179/}{410 U.S. at 211--14}). Notably, in footnote 2/4, Douglas rejected Justice Stewart's characterization of \textit{Griswold} as reviving substantive due process, insisting instead on peripheral Bill of Rights guarantees, and in footnote 2/1 he disagreed with dismissing Dr. Hallford's complaint under \textit{Younger} (\href{https://supreme.justia.com/cases/federal/us/410/179/}{410 U.S. at 209 n.2/1, 212 n.2/4}).

\subsection{Concurring Opinion (Justice Stewart)}

Justice Stewart candidly acknowledged that, despite \textit{Ferguson v. Skrupa}, 372 U.S. 726 (1963), having purported to bury economic substantive due process, \textit{Griswold v. Connecticut} could ``be rationally understood only as a holding that the Connecticut statute substantively invaded the `liberty' that is protected by the Due Process Clause of the Fourteenth Amendment'' (\href{https://supreme.justia.com/cases/federal/us/410/113/}{410 U.S. at 167--68}). Invoking Justice Harlan's \textit{Poe v. Ullman} dissent and \textit{Eisenstadt v. Baird}, Stewart concluded that personal choice in matters of family and procreation sits squarely within Fourteenth Amendment substantive liberty and that Texas's inflexible ban failed ``particularly careful scrutiny'' (\href{https://supreme.justia.com/cases/federal/us/410/113/}{410 U.S. at 169--71}).

\subsection{Dissenting Opinion (Justice White, joined by Justice Rehnquist)}

Justice White attacked both \textit{Roe} and \textit{Doe} as an illegitimate judicial coup. He argued that ``nothing in the language or history of the Constitution'' supported the majority's newly fashioned right, which elevated ``the convenience, whim, or caprice of the putative mother'' over ``the life or potential life of the fetus'' (\href{https://supreme.justia.com/cases/federal/us/410/179/}{410 U.S. at 221--22}). Calling the judgment ``an exercise of raw judicial power'' and ``an improvident and extravagant exercise of the power of judicial review,'' White contended that weighing fetal existence against maternal interests was a deeply contested moral choice that belonged to the people and state legislatures through the democratic process (\href{https://supreme.justia.com/cases/federal/us/410/179/}{410 U.S. at 222}).

\subsection{Dissenting Opinion (Justice Rehnquist)}

Justice Rehnquist raised three distinct objections. First, on standing and judicial restraint, the record failed to show that Roe was in her first trimester when she sued, so the Court violated \textit{Ashwander} avoidance principles by crafting a sweeping first-trimester rule for a hypothetical plaintiff (\href{https://supreme.justia.com/cases/federal/us/410/113/}{410 U.S. at 171--72}). Second, on the merits, a surgical operation by a licensed physician is not ``private'' in the ordinary or Fourth Amendment sense; rather, the liberty asserted was subject to deferential rational-basis review under \textit{Williamson v. Lee Optical}, which Texas's statute easily satisfied (\href{https://supreme.justia.com/cases/federal/us/410/113/}{410 U.S. at 172--73}). Comparing the majority's approach to \textit{Lochner v. New York}, Rehnquist emphasized that by 1868, when the Fourteenth Amendment was ratified, at least 36 states or territories had enacted statutes limiting abortion---including Texas's own 1857 law---proving that the drafters never intended to withdraw abortion regulation from the states (\href{https://supreme.justia.com/cases/federal/us/410/113/}{410 U.S. at 174--77 \& nn.1--2}). Third, on remedy, even under the majority's framework, the Court should have invalidated the statute only as applied rather than striking it down \textit{in toto} (\href{https://supreme.justia.com/cases/federal/us/410/113/}{410 U.S. at 177--78}).

\section{Examples: Future Applications}

Under the doctrinal framework announced in \textit{Roe v. Wade} (evaluated on its own terms as decided in 1973):
\begin{enumerate}
\item \textbf{Same-Side (Invalidated): Mandatory First-Trimester Hospitalization Law.} A state enacts a law requiring that all abortions, including those performed at eight weeks of gestation, take place in a full-service licensed hospital rather than an outpatient clinic. \textit{Why:} Under \textit{Roe}'s first-trimester rule, the state's interest in maternal health is not yet compelling because early abortion is as safe as or safer than childbirth; thus, the state may not regulate the facility or procedure during the first trimester beyond requiring a licensed physician (cf.\ \href{https://supreme.justia.com/cases/federal/us/410/179/}{\textit{Doe v. Bolton}}, 410 U.S. at 195).
\item \textbf{Same-Side (Invalidated): Post-Viability Ban Without a Maternal Health Exception.} A state prohibits all abortions after 26 weeks of pregnancy unless necessary to prevent the woman's imminent death, omitting any exception for serious threats to her physical or mental health. \textit{Why:} Even after fetal viability, when the state's interest in potential life becomes compelling, \textit{Roe} expressly requires an exception whenever abortion is necessary, in appropriate medical judgment, to preserve the mother's ``life \textit{or health}'' (\href{https://supreme.justia.com/cases/federal/us/410/113/}{410 U.S. at 164--65}).
\item \textbf{Opposite-Side (Upheld): Licensing Requirement Barring Non-Physicians.} A state prosecutes an unlicensed practitioner who performs a first-trimester abortion in a private residence. \textit{Why:} \textit{Roe} explicitly held that a state may define ``physician'' to mean only a state-licensed physician and may proscribe any abortion performed by a non-physician at any stage of pregnancy (\href{https://supreme.justia.com/cases/federal/us/410/113/}{410 U.S. at 165}).
\item \textbf{Opposite-Side (Upheld): Second-Trimester Clinic Safety and After-Care Regulations.} A state requires clinics performing abortions at 16 weeks of pregnancy to maintain sterile surgical equipment, emergency hospital-transfer agreements, and licensed nursing staff. \textit{Why:} After the end of the first trimester, the state's maternal-health interest becomes compelling, permitting regulations reasonably related to the preservation and protection of maternal health, including facility and staffing standards (\href{https://supreme.justia.com/cases/federal/us/410/113/}{410 U.S. at 163}).
\item \textbf{Fence-Sitter (Borderline): Mandatory 24-Hour Waiting Period and Informed-Consent Script in the First Trimester.} A state requires every woman seeking an abortion at nine weeks to receive state-drafted medical information and wait 24 hours before the procedure. \textit{Why:} Under strict \textit{Roe} doctrine, this is a close, contested edge case: challengers argue it directly interferes with the first-trimester physician-patient decision, which \textit{Roe} said must be ``free of interference by the State'' (\href{https://supreme.justia.com/cases/federal/us/410/113/}{410 U.S. at 163}); defenders argue that informed-consent rules apply across ordinary medical practice and merely ensure voluntary, informed consent without banning the procedure. (Historically, the Court struck down waiting periods under \textit{Roe} in \textit{City of Akron v. Akron Center for Reproductive Health} (1983) before upholding a 24-hour waiting period under \textit{Casey}'s ``undue burden'' standard in 1992.)
\end{enumerate}

\section{Critique}

\textbf{Originalist and Textualist Critiques:} Originalists and textualists have long viewed \textit{Roe} as a textbook example of unmoored substantive due process. Echoing Justice Rehnquist's dissent, critics emphasize that the text of the Constitution mentions neither privacy nor abortion, and that 36 states and territories restricted abortion when the Fourteenth Amendment was ratified in 1868 (\href{https://supreme.justia.com/cases/federal/us/410/113/}{410 U.S. at 174--76}; \href{https://fedsoc.org/fedsoc-review/what-will-settle-dobbs}{Federalist Society Review}). In \href{https://www.supremecourt.gov/opinions/21pdf/19-1392_6j37.pdf}{\textit{Dobbs v. Jackson Women's Health Organization}} (2022), Justice Alito drew on this scholarship to argue that \textit{Roe}'s historical narrative was flawed, that its trimester scheme read like a statutory code rather than constitutional interpretation, and that removing a deeply contested moral question from state legislatures short-circuited democratic self-government (\href{https://drakelawreview.org/wp-content/uploads/2024/06/kende-9.0.pdf}{Drake Law Review}).

\textbf{Progressive, Doctrinal, and Feminist Critiques:} Remarkably, some of \textit{Roe}'s sharpest doctrinal critics were political supporters of abortion access. In his landmark 1973 article \href{https://www.aei.org/wp-content/uploads/2023/07/AEIReprint015.pdf}{\textit{The Wages of Crying Wolf: A Comment on Roe v. Wade}}, 82 Yale L.J. 920 (1973), liberal scholar John Hart Ely argued that even though he favored legalized abortion as a policy matter, \textit{Roe} was ``bad because it is bad constitutional law, or rather because it is not constitutional law and gives almost no sense of an obligation to try to be'' (\href{https://en.wikipedia.org/wiki/John_Hart_Ely}{Ely Summary}). Similarly, then-Judge Ruth Bader Ginsburg argued in \href{https://scholarship.law.unc.edu/nclr/vol63/iss2/4/}{\textit{Some Thoughts on Autonomy and Equality in Relation to Roe v. Wade}}, 63 N.C. L. Rev. 375 (1985), that \textit{Roe} was weakened by two choices: first, by sweeping away nearly all state laws at once instead of invalidating Texas's extreme statute and letting a broader legislative reform movement mature; and second, by grounding the right in medicalized physician-patient ``privacy'' rather than in the Equal Protection Clause and women's equal citizenship (\href{https://www.govinfo.gov/content/pkg/GPO-CHRG-GINSBURG/pdf/GPO-CHRG-GINSBURG-3-2-1-2.pdf}{Ginsburg Analysis}; \href{https://law.yale.edu/sites/default/files/documents/pdf/Faculty/Siegel_RoeRageDemocraticConstitutionalismAndBacklash.pdf}{Post \& Siegel, \textit{Roe Rage}}). Feminist scholars have also noted that Blackmun's opinion repeatedly framed the right around the ``attending physician's'' clinical judgment rather than women's bodily autonomy (\href{https://supreme.justia.com/cases/federal/us/410/113/}{410 U.S. at 163--66}).

\textbf{Underlying Competing Values:} At bottom, \textit{Roe} reflects an intractable clash among three foundational values: (1) individual bodily autonomy, privacy, and women's equal participation in civic and economic life; (2) the moral status and state interest in protecting prenatal human life; and (3) democratic legitimacy and judicial restraint---specifically, whether unelected federal judges or elected state legislatures should resolve contested moral questions not explicitly settled by constitutional text.

\section{Key Quotations}

\begin{quote}
``This right of privacy, whether it be founded in the Fourteenth Amendment's concept of personal liberty and restrictions upon state action, as we feel it is, or, as the District Court determined, in the Ninth Amendment's reservation of rights to the people, is broad enough to encompass a woman's decision whether or not to terminate her pregnancy.''
\end{quote}
\noindent --- Justice Blackmun, Opinion of the Court, \href{https://supreme.justia.com/cases/federal/us/410/113/}{410 U.S. at 153} (\href{https://www.courtlistener.com/opinion/108713/roe-v-wade/}{CourtListener}).

\begin{quote}
``All this, together with our observation, \textit{supra}, that, throughout the major portion of the 19th century, prevailing legal abortion practices were far freer than they are today, persuades us that the word `person,' as used in the Fourteenth Amendment, does not include the unborn.''
\end{quote}
\noindent --- Justice Blackmun, Opinion of the Court, \href{https://supreme.justia.com/cases/federal/us/410/113/}{410 U.S. at 158} (\href{https://www.courtlistener.com/opinion/108713/roe-v-wade/}{CourtListener}).

\begin{quote}
``We need not resolve the difficult question of when life begins. When those trained in the respective disciplines of medicine, philosophy, and theology are unable to arrive at any consensus, the judiciary, at this point in the development of man's knowledge, is not in a position to speculate as to the answer.''
\end{quote}
\noindent --- Justice Blackmun, Opinion of the Court, \href{https://supreme.justia.com/cases/federal/us/410/113/}{410 U.S. at 159} (\href{https://www.courtlistener.com/opinion/108713/roe-v-wade/}{CourtListener}).

\begin{quote}
``As an exercise of raw judicial power, the Court perhaps has authority to do what it does today; but, in my view, its judgment is an improvident and extravagant exercise of the power of judicial review that the Constitution extends to this Court.''
\end{quote}
\noindent --- Justice White, dissenting (joined by Justice Rehnquist), \href{https://supreme.justia.com/cases/federal/us/410/179/}{410 U.S. at 222} (\href{https://www.courtlistener.com/opinion/108713/roe-v-wade/}{CourtListener}).

\begin{quote}
``The decision here to break pregnancy into three distinct terms and to outline the permissible restrictions the State may impose in each one, for example, partakes more of judicial legislation than it does of a determination of the intent of the drafters of the Fourteenth Amendment.''
\end{quote}
\noindent --- Justice Rehnquist, dissenting, \href{https://supreme.justia.com/cases/federal/us/410/113/}{410 U.S. at 174} (\href{https://www.courtlistener.com/opinion/108713/roe-v-wade/}{CourtListener}).

\end{document}
